\subsection{Ranks of small pointed models}
\label{sec:p03-infinitary-boundedness}

The fragment lemmas give small elementary submodels and elementary
unions. When a sentence permits only well-ordered interpretations,
these two facts turn elementary extensions of pointed models into
a well-founded relation on a set. Its rank gives a uniform ordinal
bound.

\begin{hbtheorem}{thm:p03-ib-rank}{proved=annex:p03.ib.1}{A bound from small pointed models}
Let \(\lambda\) be an infinite cardinal and let \(F\) be a fragment
in the sense of \cref{def:p03-w03-p03-039-fragment}, with
\(|F|\leq\lambda\) and vocabulary of size at most \(\lambda\).
Every formula has finitely many free variables. Suppose
\(\psi\in F\) is a sentence and \(\delta(x),\rho(x,y)\in F\),
and that in every model \(\mathfrak M\models\psi\), the restriction
of \(\rho^{\mathfrak M}\) well-orders
\(D^{\mathfrak M}=\{a:\mathfrak M\models\delta(a)\}\).
Then there is an ordinal \(\beta<(2^\lambda)^+\) such that all
these interpreted well-orders have order type less than \(\beta\).
Empty carriers are allowed and have type zero. The nonempty-model
fragment lemmas are applied only to models with a marked carrier
point.
\end{hbtheorem}

\begin{hbcorollary}{thm:p03-ib-individual}{proved=annex:p03.ib.2}{Individual bounds}
Every \(L_{\infty\omega}\)-sentence that makes a designated
binary relation a well-order of its field in all models has a
uniform ordinal bound on those order types. Nevertheless, for each
ordinal \(\alpha\geq2\), a sentence of \(L_{\infty\omega}\)
pins down \(\alpha\). Thus this logic is bounded in the
sentence-by-sentence sense while its well-ordering number is
\(w(L_{\infty\omega})=\infty\), as in
\cite[Ch.~II, Example~5.2.3(e), p.~62]{Ebbinghaus1985ExtendedLogics}.
\end{hbcorollary}

The next theorem applies the rank bound to the particular
renaming-and-relativization pullbacks of Barwise's framework.
Only the symbols occurring in the witnessing sentence and the
images of the carrier and order symbols are needed.
